\begin{center}
\SourceLeaf{023b}\pdfbookmark[2]{5.}{article-005}5.
\end{center}
Angulum $v$, qui pro parabola intra terminos $-180^0$ et $+180^0$, pro hyperbola intra
$-(180^0-\psi)$ et $+(180^0-\psi)$ coërcetur, pro ellipsi vero circulum integrum periodis
perpetuo renovatis percurrit, corporis moti \emph{anomaliam veram} nuncupamus. Hactenus
quidem omnes fere astronomi anomaliam veram in ellipsi non a perihelio sed ab aphelio
inchoare solebant, contra analogiam parabolae et hyperbolae, ubi aphelium non datur
adeoque a perihelio incipere oportuit: nos analogiam inter omnia sectionum conicarum
genera restituere
\SourceLeaf{024}eo minus dubitavimus, quod astronomi gallici recentissimi exemplo suo iam praeiverunt.

Ceterum expressionis $r = \frac{p}{1+e\cos v}$ formam saepius aliquantulum mutare
convenit; imprimis notentur formae sequentes:
\[
r = \frac{p}{1+e-2e\sin\frac{1}{2}v^2} = \frac{p}{1-e+2e\cos\frac{1}{2}v^2}
\]
\[
r = \frac{p}{(1+e)\cos\frac{1}{2}v^2+(1-e)\sin\frac{1}{2}v^2}\cdot
\]

In parabola itaque habemus $r = \frac{p}{2\cos\frac{1}{2}v^2}$; in hyperbola expressio
sequens imprimis est commoda $r = \frac{p\cos\psi}{2\cos\frac{1}{2}(v+\psi)\cos\frac{1}{2}(v-\psi)}\cdot$
